\section{Auditing Proofs through Counterexamples}
\label{sec:introduction}
Mathematical arguments, whether written by humans or generated by AI systems \citep{hubert2025alphaproof,feng2026firstproof}, can fail at a single intermediate step while the surrounding argument reads well. A proof may rely on an identity that fails at a boundary case, apply a lemma outside its hypotheses, or claim an invariant that breaks after several steps. Finding and diagnosing such errors is part of mathematical work, and machine-generated arguments add to what must be checked.

Delegating that review to an automated critic raises an immediate problem: a critic's objection can be wrong too. It may challenge a true statement, ignore a hypothesis inherited from an earlier paragraph, or flag a flawed step with an invalid example. Auditing a proof is also different from solving the problem afresh: a new proof of a true theorem does not locate the defect in an earlier argument for it. Useful criticism therefore needs evidence that a reader or tool can inspect directly. A counterexample provides that evidence: a concrete object that satisfies the assumptions of an assertion and violates its conclusion. In \emph{Proofs and Refutations}, \citet{lakatos1976proofs} analyzes this interplay of proofs, counterexamples, and revisions. Examining why a concrete example breaks an assertion can expose a hidden assumption, guide a repair, or reveal a pattern that extends beyond the original case.

Three questions guide this paper. \textbf{Why counterexamples?} A counterexample refutes a specific assertion with a concrete witness that readers can check, even when the theorem being proved is true. \textbf{When can we trust a counterexample?} Only after verifying that the witness satisfies the source's assumptions and contradicts its claimed conclusion, and determining whether the failure is local to one step or falsifies the main theorem. \textbf{What can we learn from a counterexample?} Analyzing why a witness works can identify missing hypotheses, guide a repair, and lead to theorems over an entire family.

We introduce \textbf{Co-PEER} (Collaborative Proof Expert Evaluation Review), a workflow for constructing and checking \emph{proofs of error}, and \textbf{BugBank}, the collection that stores them. Each \emph{objection dossier} in BugBank records the challenged source passage and its assumptions, a concrete counterexample, and a formal \emph{certificate}: a Lean-checked proof that the encoded assumptions hold and the encoded conclusion fails for the witness. Each dossier also records its revision history and a separate review of \emph{source correspondence}, whether the formal statement represents the claim in the source. Lean checks the certificate under its encoded definitions, whereas the source review decides whether the objection applies to the original argument. For our localization benchmark, we select a \emph{human-approved core} of 147 dossiers from 137 distinct sources using recorded human judgments on source interpretation, counterexample validity, and correspondence with the Lean statement (Section~\ref{sec:evaluation}).

A single counterexample can also open a broader mathematical question. Consider the 18-vertex graph $C_9[K_2]$, formed by replacing each vertex of a 9-cycle with a pair of adjacent vertices and joining each vertex of a pair to both vertices of the neighboring pairs. This graph needs five colors, yet it also admits a proper six-coloring in which every vertex already sees all five other colors among its neighbors, so no single vertex can change color while keeping the coloring proper. Asked to prove or refute a conjecture of \citet{belavadi2026recoloring}, with no access to BugBank, a model produced this \emph{frozen} six-coloring of $C_9[K_2]$. A graph is \emph{recolorable} if any two proper colorings can be joined by single-vertex changes whenever the palette exceeds the chromatic number. The conjecture states that if a prime graph and all its induced subgraphs are recolorable, then every clique blowup of the graph is recolorable as well. The base cycle $C_9$ satisfies these hypotheses, so the frozen coloring of $C_9[K_2]$ refutes the conjecture.

Studying why this coloring freezes yielded a classification for doubled odd cycles $C_n[K_2]$ of length $n\ge5$. Frozen six-colorings exist if and only if $3\mid n$, and when $3\mid n$ there are exactly $90\cdot2^n$ such colorings for labeled vertices and colors. The existence construction was known \citep{shaebani2009fall}; our analysis proves necessity and the exact count (Section~\ref{sec:mathematics}).

\paragraph{Contributions and principal findings.}
We introduce \textbf{Co-PEER} and \textbf{BugBank} to make mathematical objections inspectable, checkable, and reusable, and we demonstrate three uses of this evidence. First, from the human-approved core we construct a benchmark of 100 proof texts, each known to contain an error (\emph{certified-claim localization}). Given a numbered proof text, a critic selects up to three candidate lines; a \emph{reference line} is one of the input's accepted error lines. The strongest evaluated LLM critic selects a reference line as its first choice in 73.0\% of texts, compared with 18.0\% for the best of four simple text rules. A separate LLM-based review of saved critic responses shows why this score alone is insufficient: some responses that select a reference line do not supply a valid objection, while some that miss it do.

Second, we study how counterexamples affect error search, proof repair, and proof revision. Across 100 proofs flagged by the workflow as flawed, counterexample-first search uses 15--31\% fewer output tokens than sequential review that must also give a counterexample. The rates of objections confirmed by a reference-assisted model grader show no clear difference between the instructions; confidence intervals for the differences include zero. On 16 Lean repair tasks, models receive the identified defect and the correction to prove. Also supplying the checked counterexample raises first-attempt completions from 13 to 15 for Sonnet~5 and from 1 to 2 for GPT-OSS 120B, while Opus~5.5 completes 15 of 16 in both conditions. In a separate study on 50 paired proof-revision tasks per model using an earlier BugBank collection, retrieving stored objections did not consistently increase the number of proofs that a model grader judged correct. The broader value of retrieval during proof development therefore remains open (Section~\ref{sec:evaluation}).

Third, with the help of AI tools, we develop selected counterexamples into explanations of entire families of failures. We determine every input on which a printed integer square-root algorithm fails and prove the correction proposed in the source. Building on the known existence construction, we prove necessity and the exact count of frozen six-colorings for doubled odd cycles of length $n\ge5$. We also prove that one two-threshold coloring of the integers is the unique optimum within its family, a construction and lower bound that independent work also obtained. These follow-ups do not test whether retrieving BugBank records helps mathematical discovery (Section~\ref{sec:mathematics}).

Section~\ref{sec:evidence} defines proofs of error and their source-correspondence requirements; Section~\ref{sec:resource} describes Co-PEER and BugBank. Section~\ref{sec:evaluation} reports the human reviews, the localization benchmark, and the search, repair, and retrieval studies; Section~\ref{sec:mathematics} develops the mathematical follow-ups. Section~\ref{sec:discussion} discusses related work and limitations, and the appendices supply supporting evidence, protocols, and proofs.
